\scanpage{431}
\mspage{343}

b) $\underline S(\pi)$ satisfait à $\mathrm L(\unsure{\Phi|\mathscr I|\mathscr D})$, et pour tt $x\in\mathrm{Ob}\,\unsure{M|\mathcal M}$, on a ceci : $x\in\mathrm{Ob}\,\underline S$ \unsure{sss|ssi} la sous-catégorie str. pleine $\underline S(\pi)_{/x}$ de $M(\pi)_{/x}$ est cofinale dans $M(\pi)_{/x}$.

\textcircled{c} $\underline S(\pi)$ satisfait à $\mathrm L(\Phi)$, et pour tt $x\in\mathrm{Ob}\,M$, et \deleted{de} \deleted{\unsure{t|tt}} \unsure{représentation} $x\simeq\varinjlim_{i\in I}x_i$ qui est $\Phi$-$\pi$-adaptée (cf. déf. plus bas), on a ceci : $x\in\mathrm{Ob}\,\underline S$ sss $I_{\underline S}$ est cofinal dans $I$, où $I_{\underline S}$ est la sous-catégorie str. pleine de $I$ formée des $i\in\mathrm{Ob}\,I$ tels que $x_i\in\mathrm{Ob}\,\deleted{\unsure{I|S}}\unsure{\underline S|\underline I}$. (ou encore, $x_i\in\mathrm{Ob}\,\underline S(\pi)$).

NB On dit qu'un système inductif $(x_i)_{i\in I}$ de $M$ est $\Phi$-$\pi$-adapté si 1) $I$ équiv. à une petite catégorie et $I$ \unsure{grde|grd} devant $\pi$ 2) les $\varinjlim$ de type $\Phi$ existent dans $I$, 3) le foncteur $i\unsure{\mapsto|-}x_i$ de $I$ dans $M$ commute aux $\varinjlim$ de type $\Phi$, et enfin 4) les $x_i\in\mathrm{Ob}\,M_\pi$.

c') ((\deleted{Si} $\pi^{\pi_0}=\pi$, où $\pi_0$ est un cardinal $\leq\unsure{\pi|\mathrm r}$ utile \unsure{à|pour} $M$, \deleted{\unsure{si on a} $\unsure{I_{\alpha_0}}$ de plus $I_{\alpha_0}$ est \unsure{grde} devant $\pi_0$, et enfin} et si \deleted{$\unsure{\Phi\ni I_{\alpha_0}}$ est} \deleted{\unsure{est}} $J\in\Phi$ implique que $J$ est grd devant $\pi_0$) Comme c), avec $I$ \uline{ordonnée}.

% slanted
\margin{\uline{NB} Quand les conditions restrictives ci-contre ne sont pas satisfaites, alors a, a', b, c sont équiv. et impliquent c', mais ne sont peut-être pas \unsure{équivalentes} à c').}

\margin{\uline{NB} Il faudrait \unsure{réénoncer|réén.} \unsure{p.|pl.} la condition en \unsure{gr|gros|$\gamma$} (si \unsure{\ldots}), cf p. 359\unsure{,|-}360}

\margin{NB Il y aurait intérêt à remonter cette définition (\uline{sans} b), déjà avant 4.13.15.}

\uline{Définition} 4.13.23. Soient $M$ une catégorie accessible, $\pi$ \inserted{\unsure{est de}} un cardinal, $\Phi$ un ens. de catégories, satisfaisant les deux conditions

\unsure{a*)|a)} $\deleted{\unsure{J|I}}\unsure{I|J}\in\Phi\Longrightarrow\mathrm{card}\,\unsure{\mathrm{red}|\mathrm{rel}}\,J\leq\pi$

\deleted{\unsure{b)}}b) $\exists$ ordinal limite $\alpha_0$, tel que $I_{\alpha_0}\in\Phi$.

% slanted
\margin{Si c'est pour poser une \uline{définition}, on peut se \unsure{contenter|contenter de} \unsure{a)|de a)} \unsure{sans|pour} \unsure{b)}. $\leftarrow$ $\nwarrow$}

% archive sheet number: 430

